\abstract{Linkage disequilibrium (LD) makes causal GWAS variants indistinguishable from correlated neighbours; resolving them is the fine-mapping problem. The challenge is species-specific: humans face dense ancestry-imbalanced LD; yeast and \emph{Arabidopsis} have exceptionally long LD; crop and livestock germplasm have sparse and fragmented functional annotations that defeat the curation pipelines that work in human biobanks. Bayesian fine-mappers integrate annotations as flat per-variant priors, discarding the relational structure linking variants to tissue-specific eQTLs, pathways and protein--protein interactions. Hierarchical belief propagation (HBP) on a variant--gene--pathway factor graph matches Bayesian baselines at 5--40$\times$ speed; an annotation-adaptive complement, graph-augmented fine-mapping (GAFM), wins 27--2 against SuSiE at weak signal. The same code fine-maps Pan-UK Biobank across four ancestries (recovering \emph{LDLR}, \emph{APOE}, \emph{LPL}, \emph{GCKR}, \emph{ANGPTL3} at single-variant resolution), yeast, \emph{Arabidopsis} and rice. An intervention across 692 leads establishes that per-variant heterogeneity of the prior, not annotation density, is the operative quantity: adding a regulatory-element flag to an otherwise uniform human cache flips HBP narrower than GAFM from 0\% to 88\% of leads. These results recast multi-omics fine-mapping as a heterogeneity-curation rather than a coverage problem, and reframe post-GWAS analysis as message passing over biological structure rather than weighted regression on flattened annotations.}
